\section{A shared certificate for a finite family of powers}
\label{sec:shared-arithmetic}

The principal change is a change of quantified witnesses.  A compiler that
expands every arithmetic relation independently repeats number-theoretic
infrastructure that can be chosen once for the entire finite system.
We prove simultaneous existence before identifying any witness slots.
The resulting identifications preserve existential solvability, rather
than the set of assignments to the old auxiliary coordinates.
These reductions are new relative to the audited direct-compiler lineage;
we make no claim that every underlying observation is new to the literature.

\subsection{The classical positive-power interface}

For an integer $a\geq2$, let $\chi_a(j),\psi_a(j)$ be the natural Pell
sequences determined by
\begin{align*}
 \chi_a(0)&=1,&\psi_a(0)&=0,\\
 \chi_a(j+1)&=a\chi_a(j)+(a^2-1)\psi_a(j),&
 \psi_a(j+1)&=\chi_a(j)+a\psi_a(j).
\end{align*}
Thus $\chi_a(j)^2-(a^2-1)\psi_a(j)^2=1$.  Write
\[
 \mathcal P(a;x,y)=x^2+y^2-(ay)^2-1.
\]
We use the classical Pell characterization of exponentiation, in the
explicit form proved by \texttt{matiyasevic} and
\texttt{eq\_pow\_of\_pell} in Mathlib~\cite{mathlibpell}.
For $n,k>0$, the graph $m=n^k$ is characterized by natural witnesses
\[
 W,A,T,Z,X,Y,U,V,B,S,R
\]
with $A,B\geq2$, $k\leq Y$, $n,k\leq W$, and
\begin{align}
 \mathcal P(A;X,Y)&=\mathcal P(A;U,V)
 =\mathcal P(B;S,R)=0,\label{eq:sa-three-norms}\\
 V&>0,\qquad Y^2\mid V,\label{eq:sa-v-divides}\\
 B&\equiv1\pmod{4Y},&B&\equiv A\pmod U,\label{eq:sa-b-crt}\\
 S&\equiv X\pmod U,&R&\equiv k\pmod{4Y},\label{eq:sa-index-congruences}\\
 X+Yn&\equiv YA+m\pmod T,\label{eq:sa-power-congruence}\\
 2An&=T+n^2+1,\qquad m<T,\label{eq:sa-t-bound}\\
 \mathcal P(W+1;A,WZ)&=0.\label{eq:sa-majorant-norm}
\end{align}
The first block recovers $X=\chi_A(k)$ and $Y=\psi_A(k)$; the last
block supplies the size bound that makes the power congruence exact.
The displayed congruence moves $Yn$ to the left so that the compiler need
not evaluate a potentially negative intermediate expression.

For clarity, the classical facts needed below are
\begin{gather}
 \psi_a(j)\geq j,\qquad \chi_a(j)\geq a^j,\qquad
 \psi_a(j)\equiv j\pmod{a-1},\label{eq:sa-classical-bounds}\\
 \psi_a(j)^2\mid\psi_a\bigl(j\psi_a(j)\bigr),\qquad
 d\mid e\Longrightarrow\psi_a(d)\mid\psi_a(e).
 \label{eq:sa-classical-divisibility}
\end{gather}
Both sequences are strictly increasing in their index.  Their dependence
on $a$ is polynomial, so
congruent base parameters give congruent values at the same index.
These are the same Pell inputs used by the inherited construction.

\subsection{Orienting the congruence quotients}

A congruence $F\equiv G\pmod H$ is often compiled using two natural
quotients, $F+Hu=G+Hv$.  Canonical completeness witnesses let us orient
the four congruences in the positive-power interface and use one quotient
each.  The orientation is proved for a chosen complete family; it is not
assumed for every old certificate.

\begin{lemma}[Oriented positive-power witnesses]
\label{lem:sa-orientation}
In the completeness direction of the positive-power interface, witnesses
can be chosen satisfying
\begin{align*}
 B&=A+Uq_B,&S&=X+Uq_S,\\
 R&=k+4Yq_R,&X+Yn&=YA+m+Tq_P,
\end{align*}
with $q_B,q_S,q_R,q_P\in\N$.  They also satisfy
$B=1+4Y(q_4+1)$ for a natural $q_4$.
\end{lemma}

\begin{proof}
Choose $X=\chi_A(k)$, $Y=\psi_A(k)$ and an auxiliary even index
$M\geq2$ for which $Y^2\mid\psi_A(M)$.  Put
$U=\chi_A(M)$ and $V=\psi_A(M)$.  The standard coprimality argument
permits a natural CRT witness $B\equiv A\pmod U$,
$B\equiv1\pmod{4Y}$.  Since $U>A$, the residue $A$ lies strictly
between zero and $U$, and therefore $B\geq A$.
This proves the first orientation and $B\geq2$.

Take $S=\chi_B(k)$ and $R=\psi_B(k)$.  Joint induction in the two
recurrences shows that both coordinates are nondecreasing in the base
parameter: their initial values agree and their nonnegative recurrence
coefficients increase with the parameter.  Thus $B\geq A$ gives
$S\geq X$.  Also $R\geq k$.  The two corresponding congruences now
have nonnegative difference quotients.

For the remaining orientation, the recurrence gives
\[
 \chi_A(k)-A\psi_A(k)=-\psi_A(k-1)\qquad(k\geq1).
\]
Consequently
\[
 L_0:=X+nY-AY=n\psi_A(k)-\psi_A(k-1)>0.
\]
The power congruence says $L_0\equiv m\pmod T$, while the canonical
choice has $0<m<T$.  If $L_0=m+Tq$ with an integer $q\leq-1$, then
$L_0\leq m-T<0$, a contradiction.  Hence $q=q_P\geq0$.
Finally, $(B-1)/(4Y)$ is a positive integer, since a zero quotient
would give $B=1$; subtracting one defines $q_4\in\N$.
\end{proof}

Each oriented equality immediately implies its old congruence, so
soundness is preserved.  In the old two-quotient convention the lift is
simply $(u,v)=(0,q)$.  Two further bounds need no independent slack.
First, $Y\geq k>0$ and $B=1+4Y(q_4+1)$ imply $B\geq5$.
Second, the equations $T=m+g+1$ and $2An=T+n^2+1$ imply $A\geq2$:
if $A\leq1$, then
\[
 2An\leq2n\leq n^2+1<T+n^2+1,
\]
contradicting the defining equality.  This inference precedes the use
of any Pell theorem and does not assume that the candidate output is
already the correct power.

\subsection{Simultaneous choice of the majorant and auxiliary pair}

\begin{theorem}[One Pell infrastructure for a finite family]
\label{thm:sa-shared-pell}
Let $(n_j,k_j)$, $1\leq j\leq r$, be a nonempty finite family with
$n_j,k_j>0$.  The conjunction $m_j=n_j^{k_j}$ admits the preceding
positive-power certificates with a common $W,A,Z,U,V,B$ and a common
$q_B$.  One may impose the exact equation
\[
 W=\sum_{j=1}^r(n_j+k_j).
\]
Calls with the same exponent may additionally share their complete
Pell-index block.
\end{theorem}

\begin{proof}
We prove completeness by an explicit order of choices.  Fix the intended
outputs and put $W=\sum_j(n_j+k_j)$.  Then $W\geq n_j,k_j$ and
$W\geq2$.  Choose
\[
 A=\chi_{W+1}(W),\qquad WZ=\psi_{W+1}(W).
\]
The ordinate is divisible by $W$ by
\eqref{eq:sa-classical-bounds}, so $Z$ is natural.  The common last
norm holds, and for every $j$,
\[
 n_j^{k_j}\leq n_j^W<(W+1)^W\leq A.
\]
The classical power-bound lemma therefore gives
\[
 0<m_j<T_j:=2An_j-n_j^2-1.
\]
No property of the original choice $\max(n_j,k_j)$, beyond its two
upper bounds, was needed.

Index the distinct supplied exponents by $e_1,\ldots,e_s$ and set
\[
 X_\ell=\chi_A(e_\ell),\quad Y_\ell=\psi_A(e_\ell),\quad
 M=2\prod_{\ell=1}^s e_\ell Y_\ell,
 \quad U=\chi_A(M),\quad V=\psi_A(M).
\]
Every $e_\ell Y_\ell$ divides $M$, so
\eqref{eq:sa-classical-divisibility} gives $Y_\ell^2\mid V$ for
every $\ell$.  Also $M\geq2$ and $V>0$.
The recurrence $\psi_a(j+2)=2a\psi_a(j+1)-\psi_a(j)$ shows
$\psi_a(j)\equiv j\pmod2$, and hence $V$ is even.
The Pell norm implies $\gcd(U,V)=1$.  Thus $U$ is odd and coprime
to every $Y_\ell$, whence
\[
 \gcd\left(U,4\prod_{\ell=1}^sY_\ell\right)=1.
\]
Since $M\geq2$, strict index growth gives $U>A$.

Choose one nonnegative CRT witness
\[
 B\equiv A\pmod U,
 \qquad B\equiv1\pmod{4\prod_\ell Y_\ell}.
\]
It satisfies $B\geq A$ and all the required individual congruences
$B\equiv1\pmod{4Y_\ell}$.  Define
$S_\ell=\chi_B(e_\ell)$ and $R_\ell=\psi_B(e_\ell)$.
Lemma~\ref{lem:sa-orientation} supplies all the natural oriented
quotients.  In particular, each positive quotient $V/Y_\ell^2$ gives
the row $V=Y_\ell^2(v_\ell+1)$ with $v_\ell\in\N$.
The same index block can be supplied to every call having exponent
$e_\ell$.  The output gap and the power-congruence quotient remain
specific to its base and output.

For soundness, take any natural solution of the shared system.
The equation for $W$ gives all the required upper bounds.  The positive
output gap and the $T_j$ equation establish $A\geq2$ as above.
For each individual call, reuse the common $W,A,Z,U,V,B$ in the
classical interface.  Its own index block proves $Y\geq k>0$,
$V>0$, $Y^2\mid V$ and $B\geq2$, and its oriented rows imply the
old congruences.  The classical theorem gives $m_j=n_j^{k_j}$.
Soundness does not require an arbitrary shared auxiliary pair to arise
from the particular product index $M$ used in the completeness proof.
\end{proof}

The product defining $M$ and the CRT choice of $B$ occur only in this
proof.  They are not variable-exponent expressions or additional
operations in the final polynomial.  The final system stores the
common norm and congruence rows and the individual finite certificates.
The released implementation shares exponent blocks only for literally
identical supplied atoms; it has twenty such blocks.  It does not merge
different expressions merely because they might have equal values on
some solutions.

\subsection{Quadratic rows and positive-input specialization}

The norm $x^2-(a^2-1)y^2=1$ can be compiled with a single product
register $h=ay$ and the quadratic residual
\[
 x^2+y^2-h^2-1=0.
\]
The witness $h$ is uniquely determined and natural.  In the main norm
the same $AY$ product is already needed by the power-congruence row.
Squares that occur as terms of a quadratic residual need no separate
register.  Likewise $Y=k+y_0$ and $T=m+g+1$ are exact affine
parametrizations of the two bounds, with explicit inverse maps
$y_0=Y-k$ and $g=T-m-1$ on old solutions.

These changes apply to the positive branch.  A general exponentiation
graph must still handle $k=0$, where $m=1$, and $n=0<k$, where $m=0$.
The present source establishes positive bases and exponents at every
emitted call, so those branches are unnecessary here.  This is a property
of the full construction, not an assumption that arbitrary power inputs
are positive.  In particular the nested calls introduced below remain
positive before their candidate factorial outputs have been identified.
An unguarded norm is quadratic in the form just displayed; multiplying
it by a Boolean selector would instead be cubic, so the simplification
is not asserted for an unchanged guarded implementation.

\section{One radix for all binomial coefficients}
\label{sec:shared-radix}

Suppose a finite family of binomial graphs has natural upper arguments
$N_i$ and lower arguments $K_i$.  The independent compiler used a
separate radix $2^{N_i}+1$ for each call.  Set instead
\[
 L=\sum_i N_i,\qquad b=2^L,\qquad q=b+1.
\]
Here $b=2^L$ is one power graph, compiled by the preceding quadratic
certificate; it is not retained as an exponentiation operator.

\begin{lemma}[Common-radix digit extraction]
\label{lem:sa-common-radix}
For every $i$, the coefficient $c_i=\binom{N_i}{K_i}$ is characterized
by natural $h_i,r_i,u_i,v_i$ satisfying
\begin{align*}
 (q+1)^{N_i}&=q^{K_i}(qh_i+c_i)+r_i,\\
 q^{K_i}&=r_i+u_i+1,\\
 q&=c_i+v_i+1.
\end{align*}
The powers in these displayed semantic equations can all be replaced
by the finite positive-power graphs when their inputs are positive.
\end{lemma}

\begin{proof}
Since $L\geq N_i$, we have $q>2^{N_i}$.  The binomial coefficients
are nonnegative and their sum is $2^{N_i}$, so every coefficient in
\[
 (q+1)^{N_i}=\sum_{j=0}^{N_i}\binom{N_i}{j}q^j
\]
is smaller than $q$.  This is therefore its base-$q$ expansion without
carries.  Its digit in position $K_i$ is $\binom{N_i}{K_i}$, with the
usual value zero for $K_i>N_i$.
The first two rows perform Euclidean division by $q^{K_i}$; the last
row restricts the next remainder to the interval $[0,q)$.
Together they select exactly that digit.  Conversely Euclidean
division supplies $h_i,r_i$ and both nonnegative slacks.
\end{proof}

The digit lemma itself also handles $N_i=0$ or $K_i=0$ under the
convention $q^0=1$.  In the positive-only released construction the
upper and lower arguments are positive, so $L>0$ and all power calls
remain in their supported branch.  Equal lower arguments now use the
same denominator power $q^{K_i}$; equal upper arguments use the same
numerator power.  Literal cache equality, rather than an unproved
algebraic simplification, determines these additional identifications.

The distinction between $L$ and the Pell majorant $W$ is essential.
The latter bounds the inputs of all power calls, including bases $q$
and $q+1$.  Setting $q=2^W+1$ would create a circular size requirement.
Our order of choices is different: fix the intended high-level arithmetic
values, then their binomial upper arguments and $L$; choose $q=2^L+1$;
then all power bases and exponents are fixed, permitting the subsequent
choice of $W,A,Z,U,V,B$.

For soundness, the shared Pell system first establishes every power
graph pointwise.  The literal equation $L=\sum_iN_i$ then gives the
bound for every digit extraction.  Thus the binomial graphs are proved
before using any factorial correctness claim.  In completeness, even
when a factorial output appears in an upper argument, we first assign
its intended value and then construct these arithmetic certificates.
No upper argument refers to a private Pell witness.

\section{A self-bounding factorial certificate}
\label{sec:self-bounding-factorial}

The inherited factorial macro chose a radius using a further large
power.  Its only purpose was to make a rational approximation close
enough to the integer $n!$.  The candidate factorial output itself
provides a sufficient radius, provided both implications are proved.

\begin{theorem}[Self-bounding factorial]
\label{thm:sa-factorial}
For $n,m\in\N$, put
\[
 Q=mn^2,\qquad C=\binom{Q+n}{n},\qquad P=(Q+1)^n.
\]
Then
\[
 m=n!\quad\Longleftrightarrow\quad
 m=\left\lceil\frac{P}{C}\right\rceil.
\]
Equivalently, $m=n!$ holds precisely when natural $r,s$ satisfy
\begin{equation}
 P+C=mC+r+1,\qquad C=r+s+1.
 \label{eq:sa-factorial-rows}
\end{equation}
\end{theorem}

\begin{proof}
Always $P,C\geq1$.  The two rows are equivalent to
\[
 P=C(m-1)+(r+1),\qquad 1\leq r+1\leq C,
\]
and hence to $m-1<P/C\leq m$.  They force $m\geq1$: if $m=0$,
the first row would give $P+C=r+1\leq C$, contradicting $P\geq1$.
Conversely the ceiling inequalities give the natural witnesses
$r=P+C-mC-1$ and $s=C-r-1$.

When $n=0$, one has $Q=0$ and $P=C=1$, so both sides assert $m=1$.
Now assume $n\geq1$ and $m\geq1$, so $Q=mn^2>0$.  Write $F=n!$
and use the exact rational identity
\[
 R:=\frac{P}{C}=F\prod_{i=1}^{n}\frac{Q+1}{Q+i}.
\]
Every factor lies in $(0,1]$, giving $0<R\leq F$.
Let $d_i=(i-1)/(Q+i)$.  These numbers lie in $[0,1)$, and
\[
 \sum_{i=1}^n d_i
 \leq\frac{n(n-1)}{2Q}
 <\frac{n^2}{2Q}=\frac{1}{2m}.
\]
Induction gives the elementary inequality
$\prod_i(1-d_i)\geq1-\sum_i d_i$.  Therefore
\begin{equation}
 R>F\left(1-\frac{1}{2m}\right).
 \label{eq:sa-factorial-ratio}
\end{equation}

For soundness, suppose $m=\lceil R\rceil$.  Since $F$ is an integer
and $R\leq F$, we have $m\leq F$.  On the other hand $R\leq m$
and \eqref{eq:sa-factorial-ratio} imply
\[
 F<\frac{m}{1-1/(2m)}
   =m+\frac{m}{2m-1}\leq m+1.
\]
The first inequality is strict, including when $m=1$.  Integrality
therefore gives $F\leq m$, and so $F=m$.

For completeness, set $m=F$.  The same estimate gives
$R>F-\tfrac12>F-1$, while $R\leq F$.  Thus $\lceil R\rceil=F$,
and the displayed natural remainder witnesses exist.
\end{proof}

The proof addresses the apparent circularity directly: $m$ is initially
an arbitrary candidate; the fixed-point condition forces it to be the
factorial.  Merely replacing the radius by a smaller expression is not
safe.  For example, the choice $Q=mn$ has a false fixed point at
$n=4,m=23$, although $4!=24$.  The quadratic radius and its strict
estimate are doing mathematical work.

In the compiler, allocate the output $m$ and replace the old radius
power by the polynomial expression $mn^2$.  Keep the remaining
binomial graph, numerator power and rows
\eqref{eq:sa-factorial-rows}.  For a positive factorial input $n$,
one has $Q+n\geq n>0$ and $Q+1\geq1$ even before proving $m=n!$.
Hence the nested binomial upper argument, lower argument, numerator
base and numerator exponent are positive without using the desired
factorial conclusion.  The power and binomial graph proofs may therefore
be applied first, followed by Theorem~\ref{thm:sa-factorial}.

Finally, all binomial coefficients, ceilings, Pell sequences, divisibility
symbols and powers with variable exponents in these sections belong to
the correctness proof and the semantic description of macros.  The
released equation contains only a fixed finite list of integer quadratic
residuals.  Their sum of squares is an ordinary polynomial of degree
four with fixed integer coefficients; no such semantic operation remains
in that polynomial.
